%==============================================================
% Advanced Scientific Writing — Final Paper (PWPSP)
% Position Paper: Direction-Aware Variant Effect Prediction for
%   Human Nicotinic Acetylcholine Receptors
%
% Target venue: PSB 2027 (Pacific Symposium on Biocomputing)
% Document class: World Scientific ws-procs11x85
%
% Build:  pdflatex → pdflatex  (or: latexmk -pdf main)
%
% SKELETON — sections marked [DRAFT], [TODO], or [PLACEHOLDER]
% await experimental results, final figures, and polishing.
%==============================================================

\documentclass{ws-procs11x85}
\usepackage{ws-procs-thm}
\usepackage{booktabs}
\usepackage{graphicx}

\begin{document}

%--------------------------------------------------------------
\title{Direction-Aware Variant Effect Prediction for Human\\
Nicotinic Acetylcholine Receptors: A Position Paper}
%--------------------------------------------------------------

\author{Hari Krishna Mohan Raj\textsuperscript{*} and Karl Kirschner}

\address{Department of Computational Chemistry,\\
Hochschule Bonn--Rhein--Sieg, Sankt Augustin, Germany\\
\textsuperscript{*}E-mail: harikrishnam0711@gmail.com}

%--------------------------------------------------------------
\begin{abstract}
Missense variants in human nicotinic acetylcholine receptor (nAChR)
subunit genes are linked to congenital myasthenic syndromes, epilepsies,
and a growing list of neurological disorders. Clinical sequencing now
produces variants of uncertain significance (VUS) faster than functional
assays can characterise them. Standard computational predictors report
whether a variant is likely damaging but cannot distinguish loss of
function (LOF) from gain of function (GOF) — a distinction that directly
determines therapeutic strategy. Recent tools such as MissION and
LoGoFunc have demonstrated that direction-aware prediction is feasible
for ion channels, yet none incorporate nAChR-specific structural biology
or subunit-level context. We argue that closing this gap requires (1)~a
curated, experimentally labelled dataset of nAChR missense variants,
(2)~a feature representation that captures subunit identity and
structural context from cryo-EM data, and (3)~a benchmarked machine-learning
pipeline optimised for small, imbalanced tabular data. We present the
architecture of such a predictor, discuss the design choices and
challenges, and report preliminary evidence that engineered sequence
and structural features carry measurable signal for separating LOF
from GOF in nAChRs. We conclude with a roadmap toward a clinically
useful direction-aware VEP for this receptor family.
\end{abstract}

\keywords{Variant effect prediction; nicotinic acetylcholine receptor;
gain of function; loss of function; ion channel; machine learning;
position paper.}

% required, do-not-remove
\copyrightinfo{\copyright\ 2026 Hari Krishna Mohan Raj, Karl Kirschner.
Open Access chapter published by World Scientific Publishing Company
and distributed under the terms of the Creative Commons Attribution
Non-Commercial (CC BY-NC) 4.0 License.}

%==============================================================
\section{Introduction}\label{sec:intro}
%==============================================================

Nicotinic acetylcholine receptors (nAChRs) are pentameric ligand-gated
ion channels that mediate fast synaptic transmission at the neuromuscular
junction and throughout the central and peripheral nervous systems.
Missense variants in the seventeen human genes that encode nAChR subunits
give rise to a spectrum of disorders: congenital myasthenic syndromes
(CMS) from muscle-type receptor mutations, autosomal dominant nocturnal
frontal lobe epilepsy (ADNFLE) from mutations in \emph{CHRNA4} and
\emph{CHRNB2}, and, increasingly, neurodevelopmental and psychiatric
phenotypes linked to rare variants in additional subunits
\cite{schaaf2011, cms_review}.

The gap between variant discovery and functional understanding has
become the central bottleneck in clinical genomics. Whole-exome and
whole-genome sequencing identify tens of thousands of missense variants
per individual, but for the overwhelming majority the functional
consequence — and therefore the clinical significance — remains unknown;
they are labelled variants of uncertain significance (VUS).\footnote{The
American College of Medical Genetics and Genomics (ACMG) guidelines
explicitly recognise that computational predictors alone are insufficient
to resolve a VUS classification; functional evidence is required for
upgrade to likely pathogenic.}  For nAChR variants in particular, the
slow throughput of gold-standard electrophysiology (patch-clamp recording
of recombinant receptors) means that fewer than five hundred variants
have been experimentally characterised across all subunits, while
population databases already contain thousands of observed missense
changes in nAChR genes.

Existing computational variant effect predictors (VEPs) address one
dimension of this problem: they estimate whether a variant is
\emph{deleterious} or \emph{pathogenic}. Tools such as SIFT,
PolyPhen-2, CADD, and AlphaMissense \cite{sift,polyphen2,cadd,alphamissense}
output a continuous or categorical damage score but provide no
information on the \emph{direction} of the functional change. This is a
critical omission for ion channels and receptors, where both LOF and GOF
can be pathogenic but demand opposite therapeutic interventions. A
patient with a GOF nAChR variant causing ADNFLE may benefit from a
channel blocker, while one with an LOF variant causing a myasthenic
syndrome may require a potentiator or cholinesterase inhibitor.
Collapsing both directions into a single ``damaging'' label erases the
information needed for precision therapy.

A new wave of direction-aware predictors has recently emerged.
funNCion \cite{funncion} targets sodium and calcium channel variants;
LoGoFunc \cite{logofunc} predicts LOF/GOF/neutral genome-wide using
ensemble learning; PreMode \cite{premode} employs graph neural networks
for mode-of-action prediction; and MissION \cite{mission}, the most
recent entry, combines protein language model embeddings with structural
features to predict direction of effect across 47 ion-channel genes.
None of these tools, however, is built around the specific biology of
nAChRs: the role of subunit identity, the distinct structural
determinants at the agonist-binding interface and the channel pore,
and the conformational coupling between the extracellular and
transmembrane domains that underlies ligand-gated gating.

This position paper makes the case that direction-aware, receptor-specific
prediction for nAChRs is both necessary and feasible. We contribute:
\begin{itemlist}
\item A curated dataset of $\sim$350 experimentally characterised nAChR
  missense variants with LOF/GOF labels extracted from primary
  electrophysiology literature (Sec.~\ref{sec:data}).
\item An engineered feature representation that combines sequence
  physicochemical properties, evolutionary substitution scores, subunit
  identity encoding, and structural descriptors derived from cryo-EM
  structures (Sec.~\ref{sec:features}).
\item A benchmark of seven machine-learning classifiers under nested
  cross-validation, identifying gradient-boosted trees as the most
  suitable model class for this small, tabular dataset
  (Sec.~\ref{sec:algorithms}).
\item A critical discussion of the open challenges — data scarcity,
  the LOF/GOF spectrum, interpretability, and the absence of a shared
  benchmark — and a roadmap for the field (Sec.~\ref{sec:discussion}).
\end{itemlist}

%==============================================================
\section{Background and Related Work}\label{sec:background}
%==============================================================

\subsection{nAChR Biology and the Clinical Importance of Direction}

nAChRs are assembled as homo- or hetero-pentamers from a repertoire of
seventeen subunits ($\alpha$1–$\alpha$10, $\beta$1–$\beta$4, $\gamma$,
$\delta$, $\varepsilon$). Muscle-type receptors (embryonic:
$(\alpha1)_2\beta1\gamma\delta$; adult: $(\alpha1)_2\beta1\varepsilon\delta$)
reside at the neuromuscular junction, while neuronal subtypes
(predominantly $\alpha4\beta2$ and $\alpha7$) mediate synaptic
transmission in the brain. Each subunit contributes a large extracellular
N-terminal domain that harbours the agonist-binding site, four
transmembrane helices (M1–M4) of which M2 lines the channel pore, and
an intracellular loop that governs trafficking and modulation
\cite{nachr_structure_2021}.

The functional consequence of a missense variant depends not only on
the chemical nature of the substitution but on where in the receptor
it occurs. A substitution in the agonist-binding pocket may alter ligand
affinity (typically LOF if it reduces binding, but GOF if it increases
agonist sensitivity); a substitution in the M2 pore-lining helix may
change channel conductance or ion selectivity; a substitution in the
intracellular loop may affect receptor trafficking to the membrane
(LOF) or desensitisation kinetics (potentially GOF if it prevents
channel closure). Subunit identity further modulates these effects:
the same substitution in $\alpha4$ versus $\alpha7$ can produce opposite
phenotypes because the two subunits participate in receptors with
different stoichiometries, expression profiles, and downstream signalling.

This complexity means that a predictor that treats all nAChR variants
uniformly — without subunit context and without structural localisation
— will systematically miss the biology that determines direction of
effect. The central hypothesis of this work is that encoding subunit
identity and structural context explicitly in the feature representation
will improve LOF/GOF classification over sequence-only or
position-agnostic approaches.

\subsection{The Landscape of Variant Effect Predictors}

Table~\ref{tab:tools} places this work in the context of existing
computational tools. We distinguish three axes: (1)~whether the tool
reports pathogenicity only or direction of effect (LOF/GOF),
(2)~whether it incorporates ion-channel-specific knowledge, and
(3)~whether it models nAChR subunit identity and structure explicitly.

\begin{table}[ht]
\tbl{Comparison of computational approaches for variant effect
prediction relevant to nAChRs.}
{\begin{tabular}{@{}lcccc@{}}
\toprule
Tool / Approach & LOF/GOF & Channel-aware & nAChR-specific & Ref. \\
\midrule
SIFT                  & No  & No  & No  & \cite{sift} \\
PolyPhen-2            & No  & No  & No  & \cite{polyphen2} \\
CADD                  & No  & No  & No  & \cite{cadd} \\
AlphaMissense         & No  & No  & No  & \cite{alphamissense} \\
ESM-1b (VEP)          & No  & No  & No  & \cite{esm1b} \\
\midrule
funNCion              & Yes & Na\textsuperscript{+}/Ca\textsuperscript{2+} & No  & \cite{funncion} \\
LoGoFunc              & Yes & No          & No  & \cite{logofunc} \\
PreMode               & Yes & No          & No  & \cite{premode} \\
MissION               & Yes & 47 channels & No  & \cite{mission} \\
\midrule
\textbf{This work}    & \textbf{Yes} & \textbf{nAChR} & \textbf{Yes} & --- \\
\bottomrule
\end{tabular}}
\label{tab:tools}
\end{table}

\subsubsection{Pathogenicity-only predictors}

SIFT \cite{sift} uses sequence conservation across homologues to
flag substitutions that deviate from the evolutionary norm.
PolyPhen-2 \cite{polyphen2} combines sequence conservation with
structural features (surface accessibility, crystallographic B-factor)
trained on HumDiv and HumVar datasets. CADD \cite{cadd} integrates
over sixty genomic annotations into a single pathogenicity score using
a linear-kernel SVM. AlphaMissense \cite{alphamissense} leverages
AlphaFold's predicted structures and MSA-based conservation to achieve
state-of-the-art proteome-wide pathogenicity prediction. All four
output a continuous score where higher values indicate greater
likelihood of pathogenicity, but none distinguishes LOF from GOF.

\subsubsection{Direction-aware predictors}

funNCion \cite{funncion} was developed for sodium and calcium channel
variants associated with epilepsy and cardiac arrhythmias. It uses a
random-forest model trained on $\sim$800 functionally characterised
variants and achieves ROC $\approx$ 0.85 on held-out Na\textsubscript{v}
variants. Its feature set includes Grantham distance, BLOSUM scores,
and position within the channel protein, but it is not trained or
validated on nAChR data.

LoGoFunc \cite{logofunc} performs genome-wide LOF/GOF/neutral prediction
using an ensemble of gradient-boosted trees trained on ClinVar labels
and features spanning conservation, biochemical properties, and
protein domain annotations. It reports balanced accuracy of $\sim$0.72
for the three-class problem, but treats all genes uniformly and does
not encode receptor-specific structural context.

PreMode \cite{premode} frames mode-of-action prediction as a
graph-neural-network problem over protein structures, achieving
strong performance on curated benchmarks. However, its reliance on
high-quality structural data and its gene-family-agnostic design mean
that nAChR-specific patterns — particularly those involving subunit
interfaces — are not explicitly modelled.

MissION \cite{mission} is the closest precursor to this work. It
combines ESM-2 protein language model embeddings with structural
features (secondary structure, relative solvent accessibility,
per-residue B-factor) to predict LOF versus GOF across 47 ion-channel
genes, reporting ROC-AUC of 0.925. MissION demonstrates that structural
features complement sequence embeddings for ion-channel variant
classification. However, it aggregates data across diverse channel
families (Na\textsubscript{v}, Ca\textsubscript{v}, K\textsubscript{v},
nAChR, GABA\textsubscript{A}, etc.) and does not encode subunit
identity or the pentameric assembly logic specific to nAChRs.

\subsection{The Gap}

No existing tool simultaneously (a)~predicts direction of effect,
(b)~is trained or validated on nAChR data, and (c)~incorporates
subunit identity and pentameric structural context. This is the gap
we aim to fill. We do not claim to outperform general-purpose tools
on pathogenicity prediction; our claim is that for nAChR variants
specifically, direction matters and existing tools do not provide it.

%==============================================================
\section{Analytical Objective}\label{sec:objective}
%==============================================================

We frame the problem as supervised binary classification: given an
nAChR missense variant defined by its subunit gene, wild-type residue,
mutant residue, and residue position, predict whether its functional
effect is loss of function (LOF) or gain of function (GOF). The model
is evaluated by its ability to recover the experimentally established
direction of effect for held-out variants under stratified cross-validation.

This framing deliberately excludes:
\begin{itemlist}
\item \emph{Neutral / benign} as a separate class. Variants with no
  measurable functional effect are currently too scarce in the nAChR
  literature to constitute a trainable class. The task is therefore
  LOF-vs-GOF among functionally consequential variants, not a
  three-way classification.
\item \emph{Mixed-effect} variants. A small number of variants exhibit
  both LOF and GOF characteristics (e.g., reduced surface expression
  but increased open probability). These are flagged in the dataset
  but excluded from training until sufficient examples exist for a
  dedicated analysis.
\item \emph{Synonymous and non-coding variants}. Only missense variants
  (single amino-acid substitutions) are considered, as they constitute
  the vast majority of clinically observed coding variants and are the
  focus of existing functional assays.
\end{itemlist}

%==============================================================
\section{Experimental Data}\label{sec:data}
%==============================================================

\subsection{Data Curation}

We assembled a dataset of experimentally characterised human nAChR
missense variants by systematic mining of primary electrophysiology
literature. Search queries across PubMed, Europe PMC, and UniProt
were designed for high recall, targeting papers that report patch-clamp
or two-electrode voltage-clamp measurements of recombinant nAChRs
carrying single amino-acid substitutions. Each retrieved paper was
manually curated to extract: the subunit gene, wild-type and mutant
residues (three-letter and single-letter codes), residue position
(numbered by signal-peptide-inclusive canonical isoform), LOF or GOF
label as stated by the original authors, the experimental assay type,
and the PubMed identifier.

The current dataset contains approximately 350 variants (218 LOF,
133 GOF) across twelve of the seventeen human subunit genes
(Table~\ref{tab:subunit_counts}). The class imbalance (62\% LOF)
reflects the literature, where LOF variants are more frequently
reported — likely because LOF phenotypes (myasthenic syndromes,
severe epilepsies) are more readily ascertained clinically than
GOF phenotypes, some of which may be subclinical or protective.

\begin{table}[ht]
\tbl{Distribution of curated variants across nAChR subunit genes.
[TODO: finalise counts after curation pass.]}
{\begin{tabular}{@{}lrrr@{}}
\toprule
Subunit & LOF & GOF & Total \\
\midrule
CHRNA1  & -- & -- & -- \\
CHRNA4  & -- & -- & -- \\
CHRNA7  & -- & -- & -- \\
CHRNB2  & -- & -- & -- \\
\multicolumn{4}{c}{\dots\ \emph{[12 subunits total]}} \\
\midrule
Total   & 218 & 133 & $\sim$351 \\
\bottomrule
\end{tabular}}
\label{tab:subunit_counts}
\end{table}

[TODO: complete the subunit-level breakdown after the final curation
pass. Currently $\sim$12 subunits are represented, with the largest
counts from CHRNA1, CHRNA4, CHRNA7, CHRNB2, CHRNE, and CHRNG.]

\subsection{Label Assignment}

LOF and GOF labels follow the original authors' determination based
on one or more of: (a)~change in maximal agonist-evoked current
amplitude, (b)~shift in agonist EC\textsubscript{50}, (c)~change in
single-channel conductance or open probability, (d)~altered receptor
surface expression (for trafficking defects), or
(e)~accelerated/decelerated desensitisation kinetics. Where authors
reported multiple functional parameters, the dominant effect on
receptor function was used. Ambiguous or mixed-effect variants were
excluded.

\subsection{Limitations of the Current Dataset}

The dataset is limited in three important respects. First, with
$\sim$350 labelled examples, it is small by machine-learning standards,
constraining model complexity and raising the risk of overfitting.
Second, labels are derived from heterogeneous assays performed in
different laboratories under different conditions (oocyte vs.\
mammalian expression, different agonist concentrations, different
temperature), introducing label noise. Third, publication bias likely
means that variants with large, clear-cut functional effects are
over-represented; variants with subtle effects or effects only
observable under specific conditions may be missing entirely. These
limitations are common to all functional-genomics curation efforts
and are not specific to nAChRs; we discuss mitigation strategies in
Section~\ref{sec:discussion}.

%==============================================================
\section{Feature Engineering}\label{sec:features}
%==============================================================

We engineer approximately 50 features organised into five categories.
The design principle is that each feature should capture a dimension
of information that is relevant to the biophysical mechanism by which
a substitution alters receptor function. A summary is given in
Table~\ref{tab:features}.

\begin{table}[ht]
\tbl{Engineered feature categories. $\Delta$ denotes mutant minus
wild-type value. [TODO: finalise feature count and add per-feature
descriptions.]}
{\begin{tabular}{@{}lll@{}}
\toprule
Category & Count & Description \\
\midrule
1. Residue physicochemical
   & $\sim$10 & Hydrophobicity, charge, size, polarity, \\
   && aromaticity (WT, MT, and $\Delta$); \\
   && side-chain volume, flexibility indices \\
2. Evolutionary
   & $\sim$5  & BLOSUM62, BLOSUM80, PAM250 scores; \\
   && Grantham distance; conservation entropy \\
3. Positional
   & $\sim$3  & Normalised sequence position; \\
   && distance to key functional motifs \\
   && (C-loop, M2 helix, agonist-binding site) \\
4. Subunit identity
   & 17       & One-hot encoding of subunit gene \\
5. Structural (cryo-EM)
   & $\sim$15 & Per-residue B-factor, SASA, secondary \\
   && structure class (helix/strand/coil), \\
   && local packing density, distance to \\
   && agonist-binding site, pore radius at \\
   && residue position \\
\bottomrule
\end{tabular}}
\label{tab:features}
\end{table}

\subsection{Rationale by Category}

\paragraph{Residue physicochemical properties.}
The chemical difference between the wild-type and mutant residue —
change in hydrophobicity, charge, size — is the most direct descriptor
of a substitution's likely biophysical impact. A Trp$\to$Gly substitution
in the M2 pore-lining helix, for example, replaces a bulky hydrophobic
side chain with a proton, drastically altering both steric packing and
local electrostatics. We use established amino-acid index values from
the AAindex database \cite{aaindex}.

\paragraph{Evolutionary substitution scores.}
BLOSUM and PAM matrices quantify how frequently a given amino-acid
substitution is observed in aligned protein families, serving as a
proxy for evolutionary tolerance. A substitution with a large negative
BLOSUM62 score is rarely observed across evolution and is more likely
to be functionally consequential — but, critically, the sign of the
score does not indicate direction (LOF vs.\ GOF).

\paragraph{Positional features.}
Where a substitution occurs within the subunit sequence determines
which functional domain it affects. We encode normalised sequence
position (residue index divided by total subunit length) and,
where structural data permit, distance to functionally annotated
motifs: the C-loop (agonist binding), the M2 pore-lining helix
(channel conductance), and the M2–M3 linker (gating coupling).

\paragraph{Subunit identity.}
A 17-dimensional one-hot vector encodes which subunit gene carries the
variant. This allows the model to learn subunit-specific patterns —
for example, that substitutions in CHRNA7 (a homomeric $\alpha$7
receptor) have different typical effect directions from those in
CHRNB2 (a neuronal $\beta$ subunit that partners with $\alpha$4).

\paragraph{Structural features.}
For variants that map to a residue resolved in an experimental cryo-EM
structure, we extract per-residue B-factor (a measure of local
flexibility), solvent-accessible surface area (SASA), secondary
structure assignment (helix, strand, or coil via DSSP), and local
packing density (number of neighbouring C$\alpha$ atoms within
10~\AA). We use structures of the human $\alpha$4$\beta$2 and
$\alpha$7 nAChRs, and the Torpedo marmorata muscle-type receptor
as a high-resolution proxy for human muscle nAChRs when human
structures are unavailable at the relevant resolution.

[TODO: list PDB IDs used and the mapping protocol for variants
not resolved in any structure.]

%==============================================================
\section{Algorithms and Experimental Design}\label{sec:algorithms}
%==============================================================

\subsection{Classifier Panel}

We benchmark seven classifiers spanning a range of model complexity:

\begin{itemlist}
\item \textbf{Logistic Regression} (L2-regularised). A linear baseline
  that tests whether the decision boundary is approximately linear in
  the feature space.
\item \textbf{Support Vector Machine} (RBF kernel). A non-linear kernel
  method that can capture pairwise feature interactions.
\item \textbf{Random Forest}. An ensemble of decision trees that handles
  mixed data types and provides feature-importance estimates.
\item \textbf{XGBoost}. Gradient-boosted trees with L1/L2 regularisation,
  widely used for tabular data competitions and biomedical
  classification tasks.
\item \textbf{LightGBM}. Gradient-boosted trees with leaf-wise growth,
  often faster and more accurate than XGBoost on small datasets.
\item \textbf{CatBoost}. Gradient-boosted trees with native categorical
  feature support and ordered boosting to reduce overfitting.
\item \textbf{Multi-Layer Perceptron} (2–3 hidden layers, ReLU, dropout).
  A simple neural baseline to test whether non-linear interactions
  beyond what trees capture are present.
\end{itemlist}

\subsection{Training and Evaluation Protocol}

To obtain unbiased performance estimates on a small dataset where a
dedicated hold-out set would be too costly, we use nested
cross-validation:

\begin{arabiclist}[(3)]
\item \textbf{Outer loop:} stratified 5-fold cross-validation. In each
  fold, 80\% of variants are used for training and hyperparameter
  tuning; the remaining 20\% are held out for evaluation.
\item \textbf{Inner loop:} stratified 3-fold cross-validation on the
  training split, used by Optuna \cite{optuna} for Bayesian
  hyperparameter optimisation (50–100 trials per model).
\item \textbf{Final evaluation:} the model with the best inner-loop
  validation F1 is retrained on the full training split and evaluated
  on the held-out outer fold. Performance is reported as the mean and
  standard deviation of accuracy, F1 (macro and GOF-specific), and
  ROC-AUC across the five outer folds.
\end{arabiclist}

Stratification preserves the LOF/GOF ratio in every split. Subunit-aware
splitting (ensuring that variants from the same subunit are not
concentrated in a single fold) is used as a sensitivity analysis to
test generalisation to subunits with few labelled examples.

\subsection{Baselines}

We compare against three baselines:

\begin{itemlist}
\item \textbf{Majority-class baseline.} Always predict LOF (the majority
  class), yielding accuracy $\approx$ 0.62 and GOF F1 = 0.
\item \textbf{BLOSUM-direction heuristic.} Predict GOF if the BLOSUM62
  score is positive, LOF otherwise. This naive rule tests whether
  evolutionary conservation alone carries directional signal.
\item \textbf{Sequence-only ablation.} The full classifier panel trained
  on physicochemical, evolutionary, and positional features but
  \emph{without} subunit identity or structural features. This
  quantifies the added value of receptor-specific encoding.
\end{itemlist}

%==============================================================
\section{Results}\label{sec:results}
%==============================================================

[PLACEHOLDER — This section will be populated once the full feature
set, final dataset, and hyperparameter tuning are complete. The
current draft reports preliminary findings from an earlier iteration;
these numbers will be replaced.]

\subsection{Preliminary Observations}

Early experiments with an earlier version of the dataset
($\sim$300 variants, $\sim$40 features, 5-fold CV without nested
resampling) yielded the following indicative results:

\begin{itemlist}
\item Gradient-boosted tree methods (XGBoost, LightGBM) consistently
  outperformed linear, kernel, and neural baselines, consistent with
  expectations for small tabular data where tree-based ensembles
  handle mixed feature types and non-linear interactions without
  extensive tuning.
\item The addition of structural features (B-factor, secondary structure,
  packing density) produced a modest but consistent improvement of
  $\sim$3--5 percentage points in GOF F1 over sequence-only features,
  suggesting that structural context carries directional signal beyond
  what sequence captures.
\item The subunit one-hot encoding was among the highest-ranked features
  by permutation importance in the gradient-boosted models, confirming
  that subunit identity is predictive of effect direction.
\item The majority-class baseline (always LOF) yielded accuracy $\sim$0.62
  and GOF F1 = 0. The BLOSUM-direction heuristic performed near chance
  (F1 $\approx$ 0.40), confirming that evolutionary conservation alone
  does not resolve direction.
\end{itemlist}

[Caveat: these numbers come from an earlier pipeline iteration and
will be updated with the final nested-CV estimates. Do not cite the
exact values.]

\subsection{Planned Analyses}

Once the final dataset and feature set are frozen, we will report:

\begin{romanlist}[(vi)]
\item Mean and 95\% confidence intervals of accuracy, F1 (LOF, GOF, macro),
  and ROC-AUC for all seven classifiers under nested CV.
\item A critical-difference diagram (Demšar) comparing classifiers
  across folds.
\item Permutation feature importance for the best-performing model,
  grouped by feature category.
\item An ablation panel quantifying the marginal contribution of each
  feature category (sequence-only, +evolutionary, +positional,
  +subunit, +structural).
\item Per-subunit performance breakdown where sample size permits
  ($n \geq 10$ per subunit).
\item Confusion matrices and error analysis: which variants are
  systematically misclassified, and do they share biochemical or
  structural characteristics?
\end{romanlist}

%==============================================================
\section{Discussion}\label{sec:discussion}
%==============================================================

\subsection{Interpretation and Significance}

[PLACEHOLDER — will be written once results are finalised. Key points
to address:]

\begin{itemlist}
\item Whether the observed performance exceeds a clinically useful
  threshold and what that threshold should be for triaging VUS.
\item Which feature categories contribute most, and what this implies
  about the biophysical determinants of LOF versus GOF in nAChRs.
\item How the nAChR-specific model compares with general
  direction-aware tools (LoGoFunc, MissION) on the nAChR subset,
  once such a comparison can be made fairly.
\item The practical interpretation of model outputs: should a
  calibrated probability be reported, or a discrete LOF/GOF call
  with confidence bands?
\end{itemlist}

\subsection{Challenges and Limitations}

\subsubsection{Data Scarcity}

With $\sim$350 labelled examples, the dataset is at least an order of
magnitude smaller than typical training sets for general VEPs (CADD
trained on $\sim$15M variants; AlphaMissense used $\sim$70M variants
from population databases). Small datasets amplify three risks:
overfitting, high-variance performance estimates, and limited
generalisability to subunits or variant types not represented in the
training data. We mitigate these through nested CV for unbiased
evaluation, L1/L2 regularisation in tree-based models, and explicit
reporting of per-subunit performance where sample size permits. But
the fundamental solution — a larger dataset — requires community-wide
curation efforts and, ideally, systematic functional characterisation
of nAChR variants at scale.

\subsubsection{The LOF/GOF Dichotomy}

Biological reality is more nuanced than a binary LOF/GOF label.
Variants can produce mixed effects (reduced surface expression but
increased open probability), subunit-dependent effects (a substitution
that is GOF in $\alpha$4$\beta$2 but LOF in $\alpha$7), or
context-dependent effects (only observable under specific agonist
concentrations or in the presence of allosteric modulators). Our
binary framing is a practical simplification that enables supervised
learning on available labels, but it should not be mistaken for a
complete description of variant biology. Future iterations could
adopt a multi-label or ordinal regression framework if data volume
permits.

\subsubsection{Label Noise and Assay Heterogeneity}

Functional labels are extracted from studies spanning decades,
laboratories, and protocols. Two-electrode voltage clamp in Xenopus
oocytes (the most common assay) differs from patch clamp in mammalian
cells in ways that can affect whether a variant is called LOF or GOF,
particularly when the effect size is small. Systematic differences in
how authors define and report LOF versus GOF introduce label noise
that limits the achievable ceiling of any supervised model trained on
these data.

\subsubsection{Structural Coverage}

Not all nAChR subunits have high-resolution cryo-EM structures, and
not all residue positions within a solved structure are resolved.
Variants in disordered loops or unresolved domains cannot be assigned
structural features and must either be excluded or imputed from
homologous structures, introducing additional uncertainty.

\subsection{Future Work}

\paragraph{Near-term (within this project cycle):}
\begin{itemlist}
\item Complete the feature-ablation study to quantify each category's
  contribution.
\item Experiment with ESM-2 protein language model embeddings as an
  alternative or complementary representation to hand-engineered
  features.
\item Implement calibration (isotonic or Platt scaling) so that model
  outputs can be interpreted as confidence estimates.
\item Package the trained model as a lightweight Python module with
  a simple CLI for prospective use by nAChR researchers.
\end{itemlist}

\paragraph{Medium-term (community effort):}
\begin{itemlist}
\item Establish a shared, versioned benchmark dataset for nAChR
  direction-of-effect prediction, analogous to the ProteinGym
  benchmark for variant effect prediction, to enable fair comparison
  across methods.
\item Expand the dataset through systematic literature mining with
  LLM-assisted extraction, targeting the long tail of functionally
  characterised but not digitally curated variants.
\item Incorporate additional receptor families (GABA\textsubscript{A},
  glycine, 5-HT\textsubscript{3}) that share the pentameric
  ligand-gated ion channel architecture, testing whether a
  multi-receptor model with shared structural representations
  outperforms receptor-specific models.
\end{itemlist}

\paragraph{Longer-term:}
\begin{itemlist}
\item Integrate the predictor with clinical variant interpretation
  pipelines (e.g., as a plug-in for the ACMG/AMP framework) to
  assess real-world impact on VUS resolution rates.
\item Investigate whether direction-of-effect prediction can guide
  \emph{in silico} drug repurposing — matching a variant's predicted
  effect direction to drugs with known opposing mechanisms.
\end{itemlist}

%==============================================================
\section{Conclusion}\label{sec:conclusion}
%==============================================================

This position paper has argued that direction-aware variant effect
prediction for human nicotinic acetylcholine receptors fills a genuine
and consequential gap in the computational genomics landscape. Existing
tools either report pathogenicity without direction or predict direction
without receptor specificity; none combines directional prediction
with nAChR subunit-level modelling. We have presented the architecture
of a predictor designed to bridge this gap: a curated, experimentally
labelled dataset; an engineered feature representation that encodes
sequence properties, evolutionary constraint, subunit identity, and
cryo-EM structural context; and a rigorous benchmarking framework under
nested cross-validation with multiple baselines.

Preliminary evidence suggests that the engineered features carry
measurable signal for distinguishing LOF from GOF, that subunit
identity and structural descriptors contribute above what sequence
alone provides, and that gradient-boosted tree methods are well-suited
to this small, tabular dataset. The main obstacles — data scarcity,
the LOF/GOF dichotomy, label noise, and limited structural coverage
— are real but not insurmountable, and we have outlined a roadmap
for addressing them.

The broader ambition is not merely to build one more variant effect
predictor, but to make the case that receptor-specific, direction-aware
prediction is the right level of resolution for ion-channel variants
and that nAChRs are a well-motivated starting point. We invite the
biocomputing community to join this effort: contribute curated variant
data, benchmark new methods against a shared nAChR test set, and help
move variant interpretation from ``damaging or not?'' to ``in which
direction, and what should we do about it?''

%==============================================================
% Acknowledgments
%==============================================================
\section*{Acknowledgments}

We thank Prof.\ Youssef Ghofrani for guidance on the structure and
argumentation of this position paper, and the Advanced Scientific
Writing course at Hochschule Bonn--Rhein--Sieg for providing the
framework within which this work was developed. H.K.M.R.\ is supported
by the MSc Autonomous Systems programme at HBRS. We acknowledge the
contributors to the open-source tools used in this work: scikit-learn,
XGBoost, LightGBM, CatBoost, Optuna, Biopython, and the Protein Data
Bank.

%==============================================================
% References
%==============================================================
\begin{thebibliography}{99}

\bibitem{sift}
P.~C. Ng and S.~Henikoff,
SIFT: predicting amino acid changes that affect protein function,
{\em Nucleic Acids Res.} {\bf 31}, 3812 (2003).

\bibitem{polyphen2}
I.~A. Adzhubei {\em et al.},
A method and server for predicting damaging missense mutations,
{\em Nat. Methods} {\bf 7}, 248 (2010).

\bibitem{cadd}
M.~Kircher {\em et al.},
A general framework for estimating the relative pathogenicity of human
genetic variants,
{\em Nat. Genet.} {\bf 46}, 310 (2014).

\bibitem{alphamissense}
J.~Cheng {\em et al.},
Accurate proteome-wide missense variant effect prediction with
AlphaMissense,
{\em Science} {\bf 381}, eadg7492 (2023).

\bibitem{esm1b}
N.~Brandes {\em et al.},
Genome-wide prediction of disease variant effects with a deep protein
language model,
{\em Nat. Genet.} {\bf 55}, 1512 (2023).

\bibitem{funncion}
A.~Brunklaus {\em et al.},
Gene variant effects across sodium channelopathies predict function and
guide precision therapy,
{\em Brain} {\bf 145}, 4275 (2022).

\bibitem{logofunc}
D.~Stein {\em et al.},
Genome-wide prediction of pathogenic gain- and loss-of-function variants
from ensemble learning of a diverse feature set,
{\em Genome Med.} {\bf 15}, 103 (2023).

\bibitem{premode}
G.~Zhang {\em et al.},
Predicting the mode of action of missense variants (PreMode),
{\em bioRxiv} (2024).

\bibitem{mission}
[Author names TBD],
MissION: direction-of-effect prediction for ion-channel missense variants
via protein language model embeddings and structural features,
{\em medRxiv} (2025).

\bibitem{nachr_structure_2021}
A.~Gharpure {\em et al.},
Agonist selectivity and ion permeation in the $\alpha$4$\beta$2 nicotinic
acetylcholine receptor,
{\em Neuron} {\bf 104}, 501 (2019).

\bibitem{schaaf2011}
C.~P. Schaaf and H.~Y. Zoghbi,
Solving the autism puzzle a few pieces at a time,
{\em Neuron} {\bf 70}, 806 (2011).

\bibitem{cms_review}
A.~G. Engel {\em et al.},
Congenital myasthenic syndromes: pathogenesis, diagnosis, and treatment,
{\em Lancet Neurol.} {\bf 14}, 420 (2015).

\bibitem{aaindex}
S.~Kawashima {\em et al.},
AAindex: amino acid index database, progress report 2008,
{\em Nucleic Acids Res.} {\bf 36}, D202 (2008).

\bibitem{optuna}
T.~Akiba {\em et al.},
Optuna: a next-generation hyperparameter optimization framework,
{\em Proc. KDD} (2019).

\bibitem{enac}
R.~P. Lifton {\em et al.},
Molecular mechanisms of human hypertension,
{\em Cell} {\bf 104}, 545 (2001).

\bibitem{mohanraj2026}
H.~K. Mohan Raj {\em et al.},
A curated dataset of experimentally characterised missense variants in
human nicotinic acetylcholine receptors, (2026).

\end{thebibliography}

\end{document}
